\documentclass[12pt]{article}

% --------------------
% PAQUETES BÁSICOS
% --------------------
\usepackage[spanish,es-noshorthands]{babel}
\usepackage[utf8]{inputenc}      % Si usas pdflatex
\usepackage[T1]{fontenc}
\usepackage{lmodern}

\usepackage{amsmath, amssymb, amsthm}
\usepackage{geometry}
\usepackage{booktabs}   
\usepackage{setspace}
\usepackage{graphicx}
\usepackage{tikz}
\usepackage{tikz-cd}
\usepackage[
  colorlinks=true,
  linkcolor=blue,
  citecolor=blue,
  urlcolor=red
]{hyperref}

\newtheorem{thm}{Teorema}[section]
\newtheorem{dfn}[thm]{Definición}
\newtheorem{lem}[thm]{Lema}
\newtheorem{ej}[thm]{Ejemplo}
\newtheorem{cor}[thm]{Corolario}
\newtheorem{obs}[thm]{Observación}
\newtheorem{con}[thm]{Conjetura}
\newtheorem{prop}[thm]{Proposición}

\DeclareMathOperator{\Frm}{Frm}
\DeclareMathOperator{\Ord}{Ord}
\DeclareMathOperator{\Top}{Top}
\DeclareMathOperator{\pt}{pt}
\DeclareMathOperator{\id}{id}
\DeclareMathOperator{\tp}{tp}
\DeclareMathOperator{\fd}{fd}
% --------------------
% CONFIGURACIÓN
% --------------------
\geometry{margin=2.5cm}
\onehalfspacing

% --------------------
% DATOS DEL DOCUMENTO
% --------------------
\title{Reporte mensual de actividades}
\date{1-febrero-2026 a 28-febrero-2026}
\author{Juan Carlos Monter Cortés}

% --------------------
% INICIO DEL DOCUMENTO
% --------------------
\begin{document}

\maketitle
% \tableofcontents

\section*{Actividades realizadas}

El presente documento tiene como finalidad reportar las actividades desarrolladas durante el periodo correspondiente. 
En este periodo el trabajo se centró principalmente en el análisis estructural del intervalo de admisibilidad, su interacción 
con levantamientos de núcleos y la formulación de preguntas abiertas relacionadas con la eficiencia y el apilamiento.\\

A continuación se describen las actividades desarrolladas:

\begin{enumerate}
  \item Se elaboró realizó la escritura de este documento con el objetivo de organizar y sistematizar los resultados obtenidos hasta el momento.
  
  \item Se analizó el comportamiento de los elementos del intervalo de admisibilidad con respecto a los levantamientos de núcleos.
  
  \item Se estudió la posible implicación entre eficiencia y las nociones libres de puntos de apilado y fuertemente apilado. 
  Como parte de este análisis surgieron las preguntas abiertas presentadas al final de la Subsección \ref{N. Apilamiento}.
\end{enumerate}

\section{Preliminares}
Una \emph{retícula completa} es un conjunto $A$ dotado de un orden parcial, con supremos e ínfimos arbitrarios,
y con elementos distinguidos que fungen como neutros de dichas operaciones; usualmente se denotan por 0
el neutro del supremo y por 1 el neutro del ínfimo.

Un ejemplo importante de retícula completa es el de los \emph{marcos}, los cuales tienen la estructura $(A, \leq, \bigvee, \wedge, 0, 1)$ y
satisfacen la \emph{ley distributiva de marcos}, es decir: 
\[
a\wedge\bigvee X=\bigvee\{a\wedge x\mid x\in X\}
\]
para todo $a\in A$ y todo $X\subseteq A$. Si una retícula completa $A$ satisface lo anterior, decimos que $A\in \Frm$, donde $\Frm$ es la categoría 
cuyos objetos son los marcos.\\

Si $A, B\in \Frm$, una función $f\colon A\to B$ es un \emph{morfismo de marcos} si preserva la estructura de retícula: es monótona,
preserva 0 y 1, preserva ínfimos finitos y supremos arbitrarios..\\

El ejemplo típico de un marco es la topología de un espacio topológico. Si $S\in \Top$, entonces $\mathcal{O}S$ (la colección de abiertos de $S$) 
es un marco con estructura 
\[
(\mathcal{O}S, \subseteq, \bigcup, \cap, \emptyset, S).
\]

Además, si $S, T\in \Top$ y $f\colon S\to T$ es continua, entonces
\[
f^{-1}\colon \mathcal{O}T\to \mathcal{O}S
\] 
es un morfismo de marcos.\\

Un hecho de gran importancia es la adjunción entre $\Frm$ y $\Top$
\[\begin{tikzcd}
\Top \arrow[rr, "\mathcal{O}(\_)", bend left] & \perp & \Frm \arrow[ll, "\pt(\_)", bend left]
\end{tikzcd}
\]
donde $\mathcal{O}(-)$ es el funtor de abiertos y $\pt(\_)$ el funtor de puntos. En particular, para $A\in \Frm$, el espacio $\pt A$ tiene como topología
\[
\mathcal{O}\pt A=\{\mathcal{U}_A(x)\mid x\in A\},\quad \mathcal{U}_A(x)=\{p\in \pt A\mid x\nleq p\}.
\]
Esto induce el morfismo de marcos $\mathcal{U}_A\colon A\to \mathcal{O}\pt A$, conocido como la \emph{reflexión espacial},
el cual es suprayectivo. Además, $\mathcal{U}_A$ es un isomorfismo si y sólo si $A$ es espacial.\\

Para $S\in\Top$, el funtor $\mathcal{O}(-)$ asigna al espacio $S$ su marco de abiertos $\mathcal{O}S$.
Así, si $V\in \mathcal{O}S$, su complemento $V'$ pertenece a $\mathcal{C}S$, la colección de subconjuntos cerrados de $S$.\\

A partir de esta información básica desarrollaremos los temas siguientes, que dividimos en dos apartados.
El primero se centra exclusivamente en cuestiones internas de marcos, mientras que el segundo aborda nociones topológicas 
relacionadas.\\

Para una exposición general sobre teoría de marcos y su interacción con la teoría de espacios topológicos
pueden consultarse \cite{Johnstone82,PicadoPultr,picado2021separation}. El objetivo de lo que sigue es analizar cómo ciertas construcciones internas en $\Frm$
reflejan y, en algunos casos, generalizan fenómenos espaciales clásicos.

\section{Cuestiones en marcos}
En esta sección recopilamos resultados clásicos sobre núcleos, filtros y cocientes en la categoría $\Frm$, los cuales constituyen el 
marco teórico necesario para el desarrollo posterior. Salvo que se indique explícitamente lo contrario, los resultados sin demostración 
corresponden a hechos conocidos en la literatura. Las proposiciones y lemas acompañados de prueba forman parte del trabajo desarrollado 
en la presente investigación.

\subsection{Núcleos}
Cuando trabajamos con marcos podemos construir nuevos marcos a partir de otros ya dados. En lo que sigue, consideramos un marco fijo $A$
y diremos que este es nuestro marco principal. Haciendo uso de este marco, definimos las siguientes clases de funciones.

\begin{dfn}\label{nucleo}
Sea $A\in \Frm$. Decimos que $j\colon A\to A$ es un \emph{núcleo} sobre $A$ si para todo $a, b\in A$ se cumplen:
\begin{enumerate}
\item Si $a\leq b$, entonces $j(a)\leq j(b)$.
\item $a\leq j(a)$.
\item $j(a\wedge b)=j(a)\wedge j(b)$.
\item $j(j(a))=j(a)$.
\end{enumerate}
\end{dfn}
Si $j$ satisface las condiciones (1)-(3), pero no necesariamente (4), diremos que $j$ es un \emph{prenúcleo}. 
Denotamos por $NA$ al conjunto de todos los núcleos sobre $A$ y por $PrA$ al de todos los prenúcleos.\\

\begin{prop}\label{ejnucleos}
Sea $a, x\in A$, las aplicaciones 
\[
u_a(x)=a\vee x,\quad v_a(x)=(a\succ x), \quad w_a(x)=((x\succ a)\succ a)
\]
pertenecen a $NA$.
\end{prop}

Entre las distintas ventajas que tenemos al trabajar en $NA$ es que cada núcleo $j$ produce un \emph{cociente}. Abordemos esto con un poco más de detalle.

\begin{dfn}\label{cociente}
Sean $A, B\in \Frm$ y $f\colon A\to B$ un morfismo de marcos. Diremos que $B$ es un \emph{cociente} de $A$ si $f$ es suprayectivo.
\end{dfn}

Si $j\in NA$, el subconjunto
\[
A_j=\{x\in A\mid j(x)=x\}
\]
es un marco. Además, la aplicación
\[
j^*\colon A\to A_j, \quad j^*(x)=j(x),
\]  
es un morfismo de marcos suprayectivo. Llamamos a $j^*$ el \emph{paso al cociente} y a $A_j$ el \emph{marco cociente} asociado a $j$.\\

El conjunto $NA$ se ordena mediante al orden natural:
\[
j\leq k \iff j(x)\leq k(x)\text{ para todo }x\in A.
\]
Visto de esta manera, podemos dotar a $NA$ de más estructura y con ello verificar el siguiente resultado.

\begin{thm}\label{Ensamble}
Si $A\in Frm$, entonces $NA\in Frm$.
\end{thm}

Este resultado muestra que la construcción de $NA$ preserva la estructura de marco, permitiendo estudiar propiedades internas de $A$ a través 
de sus núcleos. Al marco $NA$ se le conoce como el \emph{ensamble} de $A$.\\ 

Podemos relacionar a un marco arbitrario con su ensamble a través 
del morfismo
\[
\eta_A\colon A\to NA,
\]
y, para $a\in A$, $\eta_A(a)=u_a$.\\

La asignación $A\mapsto NA$ define un funtor
\[
N\colon \Frm\to \Frm,
\]
lo cual será fundamental en el análisis posterior de levantamientos. Más detalles sobre $NA$ y $\eta_A$ se pueden consultar en 
\cite{simmons2006assembly,simmonsbasics,simmonscollection,Zaldivar2023}.\\

Para cada $j\in NA$, este tiene una representación por medio de los núcleos básicos
presentados en la Proposición \ref{ejnucleos}.

\begin{prop}\label{basenucleos}
Si $j\in NA$, se cumple que 
\[
\bigvee\{u_{j(a)}\wedge v_a\mid a\in A\}=\, j\, =\bigwedge\{w_a\mid a\in A\}.
\]
\end{prop}

Si tenemos $B\in \Frm$ y un morfismo de marcos $f\colon A\to B$ podemos obtener un núcleo usando el \emph{adjunto derecho} de $f$.

\begin{dfn}\label{adjunto}
Sea $f\colon A\to B$ un morfismo de marcos. Diremos que $f_*\colon B\to A$ es el adjunto derecho de $f$ (y escribimos $f\dashv f_*$) 
si para todo $a\in A$ y $b\in B$ se cumple
\[
f(a)\leq b\Longleftrightarrow a\leq f_*(b).
\]
\end{dfn}

En la categoría $\Frm$, todo morfismo de marcos admite adjunto derecho. Si $f$ y $f_*$ interactúan a través de la composición, esta cumple ciertas
propiedades especiales, dependiendo la forma en la que se compongan.

\begin{prop}\label{propadjunto}
    Para $f=f^*\dashv f_*$. Entonces:
    \begin{enumerate}
        \item $f\circ f_*\leq \id$ y $\id\leq f_*\circ f$.
        \item $f\circ f_*\circ f=f$ y $f_*\circ f\circ f_*=f_*$.
        \item Si $f$ es suprayectivo, entonces $f\circ f_*=\id$.
        \item La composición $f_*\circ f$ es un núcleo sobre $A$.   
    \end{enumerate}
\end{prop}

Sabemos que si $j\in NA$, entonces $A_j$ es un marco y $j^*$ es un morfismo de marcos. De esta manera, $j^*$ tiene adjunto derecho $j_*$.
\begin{prop}
     En la situación anterior, $j_*$, es la inclusión $\iota\colon A_j\to A$. 
\end{prop}

Si $f\colon A\to B$ es un morfismo de marcos, por la funtorialidad de $N(\_)$ tenemos el siguiente diagrama conmutativo
\[\begin{tikzcd}
	A & B \\
	NA & NB
	\arrow["f", from=1-1, to=1-2]
	\arrow["{\eta_A}"', from=1-1, to=2-1]
	\arrow["{\eta_B}", from=1-2, to=2-2]
	\arrow["{N(f)}"', from=2-1, to=2-2]
\end{tikzcd}\]
donde $N(f)(j)=\bigvee\{u_{f(j(a))}\wedge v_{f(a)}\mid a\in A\}$ para $j\in NA$.\\

Notemos que $N(f)$ puede ser evaluado en cualquier núcleo, entonces al evaluarlo en 
en $u_a$ y $v_a$ obtenemos
\[
N(f)(u_a)=u_{f(a)}\quad\mbox{ y }\quad N(f)(v_a)=v_{f(a)}.
\]

Propiedades adicionales de $N(f)$ las enunciamos a continuación.

\begin{prop}\label{levantamiento}
Si $A, B\in \Frm$ y $f\colon A\to B$ un morfismo de marcos, las siguientes afirmaciones se cumplen.
\begin{enumerate}
\item $N(f)$ es un morfismo de marcos,
\item Si $f$ es suprayectivo, entonces $N(f)$ es suprayectivo,
\item Si $f$ es epi, entonces $N(f)$ es epi,
\item Para cada $k\in NB$, $(N(f))_*(k)=f_*\circ k\circ f\in NA$,
\item $N(f)(j)\leq k\iff f\circ j\leq k\circ g$,
\item Si $f$ es suprayectivo, entonces $f\circ (N(f))_*(k)=k\circ f$.
\end{enumerate}
\end{prop}

Al núcleo $f_*\circ k\circ f$ lo llamamos el \emph{levantamiento} de $k$ a través de $f$.\\

Los resultados 

\subsection{Filtros}
Otra característica que tienen los marcos es que distintos objetos definidos dentro de un marco están en correspondencia biyectiva con objetos dentro del mismo marco 
y, en ocasiones, con objetos dentro de otras categorías. Ejemplo de esto son los filtros.

\begin{dfn}\label{filtros}
Sea $A\in \Frm$ y $F\subseteq A$. Decimos que $F$ es un filtro si:
\begin{enumerate}
\item $1\in F$,
\item si $a, b\in F$, entonces $a\wedge b\in F$,
\item si $a\in F$ y $a\leq b$, entonces $b\in F$.
\end{enumerate}
\end{dfn}

Equivalentemente, los filtros pueden describirse como secciones superiores dirigidas hacía abajo. Podemos definir diferentes tipos de filtros, mencionamos los que son más relevantes en lo que aquí se trabaja.

\begin{dfn}\label{tfiltros}
Sea $F\subseteq A$ un filtro. Decimos que:
\begin{itemize}
\item $F$ es \emph{propio} si $0\notin F$.
\item $F$ es \emph{primo} si es propio y $a\vee b\in F$ implica $a\in F$ o $b\in F$.
\item $F$ es \emph{completamente primo} si es propio y para todo $X\subset A$ tal que $\bigvee X\in F$, se cumple que $X\cap F\neq \emptyset$.
\item $F$ es \emph{abierto} si para todo subconjunto dirigido $X\subseteq A$, tal que $\bigvee X\in F$, se cumple $X\cap F\neq \emptyset$.
\item $F$ es \emph{admisible} si existe $j\in NA$ tal que 
\[
F=\{a\in A\mid j(a)=1\}.
\]
\end{itemize}
\end{dfn}

Denotamos por $A^\wedge$ a la familia de todos los filtros abiertos en $A$ y por $\text{CFil}$ a la familia de todos los filtros completamente primos. Los filtros abiertos están en correspondencia biyectiva con los subconjuntos 
\emph{compactos saturados} de $\pt A$, mientras que los filtros completamente primos corresponden a los puntos del marco.

Sea $f\colon A\to B$ un morfismo de marcos con adjunto derecho $f_*\colon B\to A$. Si $F\subseteq A$ y $G\subseteq B$ son filtros, definimos:
\begin{equation}\label{Imagenfiltros}
b\in f^*F \Longleftrightarrow f_*(b)\in F\qquad a\in f_*G \Longleftrightarrow f(a)\in G.
\end{equation}

\begin{prop}\label{fF}
Para $f=f^*\colon A\to B$ un morfismo de marcos y $G\in B^\wedge$, entonces $f_*G\in A^\wedge$.
\end{prop}

\begin{proof}
Por (\ref{Imagenfiltros}), $f_*G$ es un filtro en $A$. Necesitamos que $f_*G$ satisfaga la condición de filtro abierto. Sea $X\subseteq A$ tal que $\bigvee X\in f_*G$, con $X$ dirigido. Entonces
\[
Y=\{f(x)\mid x\in X\}
\] 
es dirigido y $f(\bigvee X)=\bigvee f[X]=\bigvee Y\in G$. Como $G$ es un filtro abierto y $f(\bigvee X)=\bigvee f[X]\in G$, existe $x\in X$ tal que $f(x)\in G$. Por lo tanto, $x\in f_*G$.
\end{proof}

Es momento de trabajar con filtros admisibles. Si $F$ es un filtro admisible, este es la preimagen del $1$ bajo algún $j\in NA$. Para hacer referencia al núcleo del cual proviene el 
filtro lo denotamos por $F=\nabla(j)$. Esta clase de núcleos permiten definir una relación de equivalencia dentro del ensamble a partir de lo siguiente. 
Para $j, k\in NA$ definimos
\[
j\sim k\iff  \nabla(j)=\nabla(k).
\]
Esta relación es de equivalencia.\\

Denotamos por $[j]$ la clase de equivalencia de $j$ y llamamos \emph{núcleo ajustado} al menor elemento de dicha clase.\\

En el siguiente lema mencionamos cuales son las principales propiedades de los núcleos ajustados y los filtros admisibles.

\begin{prop}\label{Fadmisiblesynucleosaju}
Para $A\in Frm$ y $F\subseteq A$ un filtro las siguientes se cumplen:
\begin{enumerate}
\item Si $f$ es un núcleo ajustado tal que $\nabla(j)=F$, entonces 
\[
f=\bigvee\{v_a\mid a\in F\},
\]
\item Si $F\in A^\wedge$, entonces $F=\nabla(j)$ para algún $j\in NA$.
\item $F\in A^\wedge$ si y solo si $A_j$ es un marco compacto.
\item Si $F\in A^\wedge$, $v_F$ es el núcleo ajustado asociado a $F$. Además, el bloque de $v_F$ siempre tiene un mayor elemento denotado por $w_F$ y se calcula de la siguiente manera 
\begin{equation}\label{wF}
w_F(a)=\bigwedge\{m\in M\mid a\leq m\}.
\end{equation}
donde $a\in A$ y $M=\{m\in \pt A\setminus F\mid m \text{ es máximo}\}$.
\end{enumerate}
\end{prop}

Volvemos a hacer uso del conjunto $M$ mencionado antes cuando hablemos del Teorema de Hofmann-Mislove.\\

Observemos que los dos últimos puntos de la Proposición \ref{Fadmisiblesynucleosaju} nos proporciona un intervalo 
\begin{equation}\label{Intervalo}
    [v_F, w_F]\subseteq NA.
\end{equation}    
A este intervalo lo llamamos \emph{intervalo de admisibilidad} asociado al filtro $F$. Este intervalo codifica todos los núcleos que producen el mismo cociente compacto.\\

Una pregunta natural bajo este contexto sería el saber cuándo
$A_j$ también es cerrado. Para responderlo, necesitamos introducir un poco más de información.\\

Con el fin de obtener el núcleo ajustado asociado a un filtro abierto, consideramos la iteración transfinita del prenúcleo correspondiente.\\

Consideremos $F\in A^\wedge$ y definimos el operador
\begin{equation}\label{Prenucleo}
f_{F}=\dot{\bigvee}\{v_a\mid a\in F\}\in PrA
\end{equation}
entonces  tenemos la sucesión creciente de prenúcleos
\begin{equation}\label{Sucesionf}
f_F^0=\id,\quad f_F^{\alpha +1}=f_F\circ f_F^\alpha,\quad f_F^\lambda=\bigvee\{f_F^\alpha\mid \alpha<\lambda\}
\end{equation}
donde $\alpha$ es un ordinal y $\lambda$ un ordinal límite.\\

De esta manera se cumple que $v_F=f_F^\infty$.\\

Evaluando cada elemento de la sucesión anterior en $0$ obtenemos la sucesión
\begin{equation}\label{Sucesiond}
d(0)=0,\quad d(\alpha + 1)=f_F(d(\alpha)),\quad d(\lambda)=\bigvee\{d(\alpha)\mid \alpha<\lambda\}.
\end{equation}
y $v_F(0)=d(\infty)=d$. Más adelante haremos uso de estos elementos para definir una clase particular de marcos.\\

Antes de concluir esta subsección es conveniente relacionar todo lo anterior con una cuestión espacial.\\

En \cite[Lema 8.9 y Corolario 8.10]{simmons2004vietoris} muestran que el diagrama
\[\begin{tikzcd}
	A & A \\
	{\mathcal{O}S} & {\mathcal{O}S}
	\arrow["{f^\infty}", from=1-1, to=1-2]
	\arrow["{U_A}"', from=1-1, to=2-1]
	\arrow["{U_A}", from=1-2, to=2-2]
	\arrow["{F^\infty}"', from=2-1, to=2-2]
\end{tikzcd}\]
conmuta laxamente, es decir, $U_A\circ f^\infty\leq F^\infty \circ U_A$. En este diagrama $f^\infty$ y $F^\infty$ representan los
núcleos ajustados asociados a los filtros $F\in A^\wedge$ y $\nabla\in \mathcal{O}S^\wedge$, respectivamente.\\

Aquí probamos algo más general, ya que consideremos $j\in NA$, $F\in A_j^\wedge$ y $j_*F\in A^\wedge$ para producir el diagrama
\[\begin{tikzcd}
	A & A \\
	{A_j} & {A_j}
	\arrow["{f^\infty}", from=1-1, to=1-2]
	\arrow["j^*"', from=1-1, to=2-1]
	\arrow["j^*", from=1-2, to=2-2]
	\arrow["{\hat{f}^\infty}"', from=2-1, to=2-2]
\end{tikzcd}\]
donde $\hat{f}^\infty$ y $f^\infty$ son los núcleos ajustados asociados a los filtros $j_*F$ y $F$, respectivamente.\\

El siguiente lema es parte del análisis desarrollado en este trabajo y constituye una herramienta clave en la construcción de los diagramas conmutativos
porteriores.

\begin{lem}\label{f1f}
    Para $j=j^*$, $f$ y $\hat{f}$ como antes se cumple que $j\circ \hat{f}\leq f\circ j$.
\end{lem}
\begin{proof}
    Recordemos que, por (\ref{Imagenfiltros})
    \[
    \hat{f}=\dot{\bigvee}\{v_y\mid y\in j_*F\} \qquad f=\dot{\bigvee}\{v_{j(y)}\mid j(y)\in F\}. 
    \]
entonces para $a\in A$ se cumple
\[
v_y(a)=(y\succ a)\leq \hat{f}(a)\leq j(\hat{f}(a)).
\]

También, para todo $a, y\in A$, $(y\succ a)\wedge y=y\wedge a$ y
\[
\begin{split}
j((y\succ a)\wedge y)\leq j(a) & \Leftrightarrow j(y\succ a)\wedge j(y)\leq j(a)\\
& \Leftrightarrow j(y\succ a)\leq (j(y)\succ j(a)).
\end{split}
\]

Así 
\[
v_y(a)\leq j(\hat{f}(a))\leq (j(y)\succ j(a))=v_{j(y)}(j(a))\leq f(j(a)).
\]

Por lo tanto $j\circ \hat{f}\leq f\circ j$.
\end{proof}

Ahora probaremos lo anterior, pero para cualquier ordinal.

\begin{cor}\label{finftyf}
    Para $j$, $f$ y $\hat{f}$ como antes, se cumple que $j\circ \hat{f}^\alpha\leq f^\alpha\circ j$
\end{cor}

\begin{proof}
    Para $\alpha\in \Ord$ verificaremos que $j\circ \hat{f}^\alpha\leq f^\alpha\circ j$. Lo haremos con inducción transfinita.\\

    Si $\alpha=0$, es trivial.\\

    Para el paso de inducción, supongamos que para $\alpha$ se cumple. Luego
    \[
    j\circ \hat{f}^{\alpha+1}=j\circ \hat{f}\circ \hat{f}^{\alpha}\leq  f\circ j\circ \hat{f}^\alpha\leq f\circ f^\alpha\circ j=f^{\alpha+1}\circ j,
    \]
    donde la primera desigualdad es el Lema \ref{f1f} y la segunda es la hipótesis de inducción.\\

    Si $\lambda$ es un ordinal límite, entonces
	\[
	\hat{f}^\lambda=\bigvee\{\hat{f}^\alpha\mid \alpha<\lambda\}, \quad f^\lambda=\bigvee\{f^\alpha\mid \alpha<\lambda\} 
	\]
	y
	\[
		j\circ \hat{f}^\lambda=j\circ \bigvee_{\alpha<\lambda}\hat{f}^\alpha\leq\bigvee_{\alpha<\lambda}j\circ \hat{f}^\alpha.
	\]
	Por la hipótesis de inducción tenemos que 
	\[
	j\circ \hat{f}^\alpha\leq f^\alpha\circ j\Rightarrow \bigvee_{\alpha<\lambda}j\circ \hat{f}^\alpha\leq \bigvee_{\alpha<\lambda} f^\alpha\circ j.
	\]
	Por lo tanto $j\circ \hat{f}^\lambda\leq f^\lambda\circ j$.
\end{proof}

\begin{obs}\label{j monotono}
El Lema \ref{f1f} y el Corolario \ref{finftyf} también se cumplen si $j$ solo es monótono y preserva ínfimos finitos.
\end{obs}

Recordando que $v_F=f^\infty$ y $v_{j_*F}=\hat{f}^\infty$, el Corolario \ref{finftyf} nos dice que $j\circ v_{j_*F}\leq v_F\circ j$.\\

El objetivo ahora es comparar los cocientes compactos inducidos por $F$ y por $j_*F$, y construir un morfismo entre ellos que haga conmutar el diagrama correspondiente.\\

Si consideramos $j, v_F$ y $v_{j_*F}$ tenemos el siguiente diagrama
\begin{equation}\label{diagrama H}
    \begin{tikzcd}
	A && {A_{j_*F}} \\
	\\
	{A_j} && {A_F}
	\arrow["{(v_{j_*F})^*}", shift left, from=1-1, to=1-3]
	\arrow["j"', from=1-1, to=3-1]
	\arrow["H"', from=1-1, to=3-3]
	\arrow["{(v_{j_*F})_*}", shift left, from=1-3, to=1-1]
	\arrow["{(v_F)^*}"', shift right, from=3-1, to=3-3]
	\arrow["{(v_F)_*}"', shift right, from=3-3, to=3-1]
\end{tikzcd}
\end{equation}
donde $A_F$ y $A_{j_*F}$ son los cocientes compactos producidos por $v_F$ y $v_{j_*F}$, respectivamente.\\

Definimos el morfismo 
\[
H\colon A\to A_F, \qquad H=v_F\circ j.
\]

Queremos definir un morfismos de marcos entre los cocientes compactos $A_{j_*F}$ y $A_F$ de tal manera que 
el diagrama (\ref{diagrama H}) sea un cuadrado conmutativo.\\

Definimos
\[
h\colon A_{j_*F}\to A_j, \qquad h=H|_{A_{j_*F}},
\]
es decir, $h(x)=H(x)$ para todo $x\in A_{j_*F}$.\\

Notemos que, por el Corolario \ref{finftyf}, se cumple que $H(A_{j_*F})\subseteq A_F$, por lo que la restricción anterior 
está bien definida. Así, si $h=H_{\mid{A_{j_*F}}}$, el diagrama 
\[\begin{tikzcd}
	A & {A_{j_*F}} \\
	{A_j} & {A_F}
	\arrow["{v_{j_*F}}", from=1-1, to=1-2]
	\arrow["j"', from=1-1, to=2-1]
	\arrow["h", from=1-2, to=2-2]
	\arrow["{v_F}"', from=2-1, to=2-2]
\end{tikzcd}\]
conmuta.\\

Debemos verificar que $h$ es un morfismo de marcos. Por la definición de $h$, este es un $\wedge$-morfismo. Restaría comprobar que es un $\bigvee$-morfismo.\\

Los supremos en $A_{j_*F}$ y $A_F$ son calculados de manera diferente. De esta manera, denotamos por $\hat{\bigvee}$ al supremo en $A_{j_*F}$ y por $\tilde{\bigvee}$ al supremo en $A_F$. Por lo tanto
\[
\hat{\bigvee}=v_{j_*F}\circ \bigvee\quad\mbox{ y }\quad \tilde{\bigvee}=v_{F}\circ \bigvee,
\] 
es decir, para $X\subseteq A$, $Y\subseteq A_j$,
\[
	\hat{\bigvee}X=v_{j_*F}(\bigvee X)\quad\mbox{ y }\quad \tilde{\bigvee}Y=v_{F}(\bigvee Y).
\]

Como $H$ es un morfismo de marcos, $H\circ \bigvee=\tilde{\bigvee}\circ H$. Necesitamos algo similar para el morfismo $h$.

\begin{lem}\label{bigvee g}
$h\circ \hat{\bigvee}=\tilde{\bigvee}\circ h$.
\end{lem}

\begin{proof}
Basta probar que 
\[
h\circ \hat{\bigvee}\leq \tilde{\bigvee}\circ h
\]
ya que la otra desigualdad se obtiene por definición de los supremos en los cocientes.\\

Notemos que
\[
h\circ \hat{\bigvee}=H\circ v_{j_*F}\circ \bigvee=v_F\circ j\circ v_{j_*F}\circ \bigvee\leq v_F\circ v_F\circ j\circ \bigvee
\]
donde la desigualdad es el Corolario \ref{finftyf}. Además, por la idempotencia de $v_F$ 
\[
h\circ\hat{\bigvee}\leq v_F\circ j\circ \bigvee =H\circ \bigvee=\tilde{\bigvee}\circ H=\tilde{\bigvee}\circ h.
\]
Por lo tanto $h\circ\hat{\bigvee}=\tilde{\bigvee}\circ h$.
\end{proof}

La construcción anterior culmina en el siguiente resultado, que establece la conmutatividad del diagrama inducido por los cocientes compactos.
\begin{prop}\label{VFsquare}
El diagrama
\begin{equation}\label{Cuadrado}
    \begin{tikzcd}
	A & {A_{j_*F}} \\
	{A_j} & {A_F}
	\arrow["{v_{j_*F}}", from=1-1, to=1-2]
	\arrow["j"', from=1-1, to=2-1]
	\arrow["h", from=1-2, to=2-2]
	\arrow["{v_F}"', from=2-1, to=2-2]
\end{tikzcd}\end{equation}
es conmutativo.
\end{prop}

Este resultado muestra que la construcción del núcleo ajustado es compatible con el paso al cociente, lo cual será fundamental en el análisis de 
levantamientos en la Sección \ref{Interacciones}.

\subsection{Axiomas de separación}
Por la relación que existe entre las categorías $\Frm$ y $\Top$, es posible definir dentro de un marco distintos tipos de axiomas de separación.
Recordamos brevemente algunos de los axiomas que serán utilizadas más adelante. Parte de lo que se incluye en esta subsección es información presentada en \cite{picado2021separation}.\\

Para poder realizar una clasificación sobre las propiedades que se mencionaran, introducimos la siguiente definción.

\begin{dfn}\label{PropFrm}
Sea $\mathcal{P}$ una propiedad definida en marcos. Decimos que la propiedad $\mathcal{P}$ es:
\begin{itemize}
\item \emph{hereditaria} si para todo marco $A$ que satisface $\mathcal{P}$ y todo $j\in NA$, el cociente $A_j$ satisface $\mathcal{P}$.
\item \emph{coproductiva} si para cada familia $\{A_i\}_{i\in I}$ de marcos que satisfacen $\mathcal{P}$, entonces $\bigoplus_{i\in I}A_i$ satisface $\mathcal{P}$, donde $\bigoplus$ representa el coproducto en $\Frm$.
\item \emph{conservativa} si existe una propiedad topológica $\mathcal{P}'$, definida sobre espacios $T_0$, tal que para todo $S\in \Top_0$,
\[
\mathcal{O}S\text{ satisface }\mathcal{P}\text{ si y solo si }S\text{ satisface }\mathcal{P}'.
\]
\end{itemize}
\end{dfn}

Consideramos $A\in Frm$ y $S=\pt A$. Sabemos que $p\in \pt A$ si y solo si $p\neq 1$ y
\begin{equation}\label{primo}
a\wedge b\leq p \Rightarrow a\leq p \; \text{o}\; b\leq p,
\end{equation}
para todo $a, b\in A$. A los elementos de $\pt A$ también podemos llamarlos \emph{primos}.\\

De manera adicional, decimos que $m\in A$ es \emph{máximo} si 
\[
a\leq m\Rightarrow a=m \text{ o } a=1
\]
para todo $a\in A$. Podemos verificar que si $m$ es máximo, entonces $m\in \pt A$. Veamos cuando el reciproco es cierto.

\begin{dfn}\label{FrmT1}
Un marco $A$ es $T_1$ si para cada $p\in \pt A$, p es máximo.
\end{dfn}

\begin{dfn}\label{Bpordebajo}
Sean $A\in \Frm$ y $a, b\in A$. Decimos que $a$ está \emph{bastante por debajo} de $b$, y lo denotamos por $a\prec b$, 
\[
\exists x\in A\text{ tal que }a\wedge c=0 \text{ y }b\vee c=1.
\]
\end{dfn}

Cuando trabajamos en marcos se cumple que para cada $a\in A$ existe $\neg a$ tal que $a\wedge \neg a=0$, la condición de la definición anterior 
se puede reescribir de la siguiente manera
\[
a\prec b\iff \neg a\vee b=1.
\]

\begin{dfn}\label{reg}
Decimos que $A\in \Frm$ es \emph{regular} (o que satisface $\mathbf{(reg)}$) si para todo $a\in A$ se cumple que
\[
a=\bigvee\{x\in A\mid x\prec a\}.
\]
\end{dfn}

Equivalentemente, podemos decir que $A$ es regular si para cada $a, b\in A$, con $a\nleq b$, existen $x, y\in A$
tales que 
\[
a\vee x=q, \quad y\nleq b, \quad x\wedge y=0.
\]

\begin{dfn}\label{aju}
Decimos que $A$ es \emph{ajustado} (o que satisface $\mathbf{(aju)})$ si para cada $a, b\in A$, con $a\nleq b$, existen $x, y\in A$
tales que 
\[
a\vee x=q, \quad y\nleq b, \quad x\wedge y\leq b.
\]
\end{dfn}

La teoría de separación en marcos es más rica que la teoría en espacios topológicos. Esto se debe a que en marcos existe más de una noción para decir
que un marco es Hausdorff.

\begin{dfn}\label{H}
Decimos que $A\in \Frm$ es \emph{Hausdorff} (o que satisface $\mathbf{(H)}$) si pata todo $a, b\in A$ tales que $1\neq a\nleq b$ existe $c\in A$ tal que
\[
c\nleq a, \qquad \neg c\nleq b.
\]
\end{dfn}

\begin{dfn}\label{fH}
Decimos que $A\in \Frm$ es \emph{fuertemente Hausdorff} (o que satisface $\mathbf{(fH)}$) si el sublocal diagonal es cerrado.
\end{dfn}

De igual manera que en el caso espacial, los axiomas en $\Frm$ se relacionan de manera jerárquica. Los siguientes diagramas muestran
como se relacionan todos los que hemos mencionado hasta este momento

\[\begin{tikzcd}
	{\mathbf{(reg)}} & {\mathbf{(fH)}} & {\mathbf{(H)}} & {T_1} \\
	& {\mathbf{(aju)}}
	\arrow[Rightarrow, from=1-1, to=1-2]
	\arrow[Rightarrow, from=1-1, to=2-2]
	\arrow[Rightarrow, from=1-2, to=1-3]
	\arrow[Rightarrow, from=1-3, to=1-4]
	\arrow[Rightarrow, from=2-2, to=1-4]
\end{tikzcd}\]

La siguiente tabla resume el comportamiento estructural de los axiomas de separación antes mencionados.
\begin{table}[htbp]
\centering
\begin{tabular}{lccccc}
\toprule
Propiedad & $T_1$ & (reg) & (H) & (fH) & (aju) \\
\midrule
Hereditaria & $\checkmark$ & $\checkmark$ & $\checkmark$ & $\checkmark$ & $\checkmark$ \\
Coproductiva & $\checkmark$ & $\checkmark$ & $\checkmark$ & $\checkmark$ &  \\
Conservativa & $\checkmark$ & $\checkmark$ & $\checkmark$ &  &  \\
\bottomrule
\end{tabular}
\caption{Comportamiento estructural de los axiomas de separación utilizados en este trabajo.}
\end{table}

\subsection{El marco de parches y marcos eficientes}
Parte de los resultados presentados en esta subsección provienen del trabajo de Rosemary A. Sexton (ver \cite{sexton2003point,sexton2006ordinal,sexton2006point}). 
Otros resultados forman parte del análisis desarrollado en esta investigación.\\

El \emph{marco de parches} de $A$, denotado por $PA$, es el submarco de $NA$ generado por supremos de los núcleos de la forma
\[
\{u_a\wedge v_F\mid a\in A, F\in A^\wedge\}. 
\]
Esto lo podemos representar a través del diagrama
\[\begin{tikzcd}
A \arrow[r] \arrow[rr, "\eta_A", bend left] & PA \arrow[r, hook] & NA
\end{tikzcd}\]

Decimos que $A$ es \emph{parche trivial} si el encaje $i\colon A\to PA$ es un isomorfimso.\\

Sexton define los marcos \emph{eficientes} y estudia la trivialidad del parche por medio de la eficiencia.

\begin{dfn}\label{eficientes}
Sea $A\in \Frm$ y $F\in A^\wedge$. Para $\alpha\in \Ord$:
\begin{enumerate}
\item $F$ es $\alpha$-eficiente si para todo $x\in F$ se cumple que $d\vee x=1$, donde $d=d(\alpha)$.
\item $A$ es $\alpha$-eficiente si cada $F\in A^\wedge$ es $\alpha$-eficiente.
\item $A$ es eficiente si $A$ es $\alpha$-eficiente para algún $\alpha\in \Ord$.
\end{enumerate}
\end{dfn}

Los marcos eficientes nos dicen que:
\begin{enumerate}
    \item La eficiencia constituye una jerarquía descendente de propiedades indexadas por ordinales. En otras palabras,
    \[
    1\text{-eficiente}\Rightarrow 2\text{-eficiente}\Rightarrow\dots\Rightarrow \alpha\text{-eficiente}\Rightarrow \dots
    \]
    \item Por construcción, se cumple que $u_d=v_F$, es decir, la eficiencia proporciona el menor cociente compacto cerrado.
    \item Sexton demuestra que un marco es eficiente si y solo si es parche trivial.
\end{enumerate}

Uno de nuestros objetivos es comprender el comportamiento de los marcos eficientes y realizar un análisis interno de estos dentro de la categoría de marcos.
Parte de la información que existe hasta el momento, relaciona la eficiencia principalmente con propiedades en $\Top$.\\

La siguiente proposición resume algunas relaciones conocidas entre eficiencia y axiomas de separación.
\begin{prop}\label{Propiedadeseficientes}
$\mbox{}$
\begin{enumerate}
\item Si $A\in \Frm$ es un marco arbitrario, entonces se cumple que
\[\begin{tikzcd}
	{\mathbf{(reg)}} & {\mathbf{(fH)}} & {\text{Eficientcia}} & {T_1} \\
	& {\mathbf{(aju)}}
	\arrow[Rightarrow, from=1-1, to=1-2]
	\arrow[Rightarrow, from=1-1, to=2-2]
	\arrow[Rightarrow, from=1-2, to=1-3]
	\arrow[Rightarrow, from=1-3, to=1-4]
	\arrow[Rightarrow, from=2-2, to=1-3]
\end{tikzcd}\]
donde $\mathbf{(reg)}$, $\mathbf{(fH)}$, $\mathbf{(aju)}$ y $T_1$ son propiedades de separación en $\Frm$.
\item Si $A=\mathcal{O}S$ para $S\in \Top$, entonces se cumple que
\begin{enumerate}
    \item $\mathcal{O}S$ es 1-eficiente si y solo si $S$ es $T_2$
    \item Para $S$ un espacio $T_0$, $S$ es eficiente ($\mathcal{O}S$ es eficiente) si y solo si $S$ es empaquetado y apilado.
\end{enumerate}
\end{enumerate}
\end{prop}
Como podemos darnos cuenta, en el caso espacial, la eficiencia está caracterizada por algunas propiedades topológicas. Estas serán abordadas en la siguiente sección.\\

Estudiamos ahora el comportamiento interno de la eficiencia respecto a construcciones categóricas. 
\begin{prop}\label{tidyquout}
    Si $A$ es un marco eficiente, entonces $A_j$ es un marco eficiente.
\end{prop}

\begin{proof}
Observemos primero que $F\subseteq j_*F$. Como $A$ es eficiente y $F\in A^\wedge$, se cumple que 
\[
x\in F\Rightarrow \hat{d}\vee x=1,
\]
donde $\hat{d}=d(\alpha)=f^\alpha(0)$.\\

Si $\hat{d}\leq d$, entonces $d\vee x=1$, para $d=d(\alpha)=f^\alpha(j(0))$.\\

Así, por el Corolario \ref{finftyf}
\[
\hat{d}=\hat{d}(\alpha)\leq j(\hat{d}(\alpha))=j(\hat{f}^\alpha(0))\leq f^\alpha(j(0))=d(\alpha)=d.
\]

Por lo tanto, si $x\in F$, entonces $d\vee x=1$ y $A_j$ es eficiente.
\end{proof}

Analizamos ahora la estabilidad de la eficiencia bajo coproductos binarios.
\begin{prop}\label{tidyplus}
Si $A_1, A_2\in \Frm$ son marcos eficientes, entonces $A_1\oplus A_2$ es un marco eficiente.
\end{prop}

\begin{proof}
Consideremos $A_1, A_2\in \Frm$ $\alpha$-eficientes. Por la construcción de $A_1\oplus A_2$ tenemos el siguiente diagrama
\[
\begin{tikzcd}
A_i \arrow[r, "\alpha_i"] \arrow[rrr, "\iota_i", bend left] & A_1\times A_2  \arrow[r, "\downarrow(\_)"] & \mathcal{L}(A_1\times A_2) \arrow[r, "\sim"] & A_1\oplus A_2 \arrow[lll, "(\iota_i)_*", bend left]
\end{tikzcd}
\]
donde $\alpha_1(a)=(a, 1)$ y $\alpha_2(b)=(1, b)$ para $a\in A_1$ y $b\in A_2$. Además, como $\iota_i\in \Frm$, $(\iota_i)_*$ es su adjunto derecho.\\

Consideremos $F\in (A_1\oplus A_2)^\wedge$ y, por la Proposición \ref{fF}, $(\iota_i)_*[F]\in A_i^\wedge$. Como $A_i$ es eficiente, para $x_i\in (\iota_i)_*F$ se cumple que $d_i\vee x_i=1$, donde 
\[
d_i=d_i(\alpha)=f_i^\alpha(0), \mbox{ con } f_i=\dot{\bigvee}\{v_y\mid y\in (\iota_i)_*F\}
\]
Por (\ref{Imagenfiltros}) se tiene que si $x_i\in (\iota_i)_*F$, entonces $\iota_i(x_i)\in F$, es decir.
\[
x_1\oplus 1, 1\oplus x_2\in F
\]
y al ser $F$ un filtro, se cumple que $x_1\oplus 1\wedge 1\oplus x_2=x_1\oplus x_2\in F$.\\

Además, para cada $i=1, 2$ tenemos el siguiente diagrama

\[\begin{tikzcd}
	{A_i} && {A_i} \\
	{A_1\oplus A_2} && {A_1\oplus A_2}
	\arrow["{f_i}", from=1-1, to=1-3]
	\arrow["{\iota_i}"', from=1-1, to=2-1]
	\arrow["{\iota_i}", from=1-3, to=2-3]
	\arrow["f_F"', from=2-1, to=2-3]
\end{tikzcd}\]
donde $f_F=\dot{\bigvee}\{v_{\iota_i(y)}\mid \iota_i(y)\in F\}$. Veamos que $\iota_i\circ f_i\leq f_F\circ \iota_i$.\\
    
Consideremos $a\in A_i$, entonces $v_y(a)\leq f_i(a)$ y $\iota_i(v_y(a))\leq \iota_i(f_i(a))$. Además,
\[
\begin{split}
\iota_i((y\succ a)\wedge y)=\iota_i(y\wedge a)\leq \iota_i(a)& \Leftrightarrow \iota_i(y\succ a)\wedge \iota_i(y)\leq \iota_i(a)\\
& \Leftrightarrow \iota_i(y\succ a)\leq (\iota_i(y)\succ \iota_i(a)).
\end{split}
\]
Luego 
\[
\iota_i(f_i(a))\leq (\iota_1(y)\succ \iota_i(a))=v_{\iota_i(y)}(\iota_i(a))\leq f_F(\iota_i(a)).
\]
Por lo tanto $\iota_i\circ f_i\leq f_F\circ \iota_i$. Verifiquemos que $\iota_i\circ f_i^\alpha\leq f_F^\alpha\circ \iota_i$ también es cierto. Por inducción transfinita 
\begin{itemize}
\item Si $\alpha=0$, es trivial.
\item Para el paso de inducción, supongamos que para $\alpha$ se cumple. Luego
\[
\iota_i\circ f_i^{\alpha+1}=\iota_i\circ f_i\circ f_i^{\alpha}\leq  f_F\circ \iota_i\circ f_i^\alpha\leq f_F\circ f^\alpha_F\circ \iota_i=f^{\alpha+1}_F\circ \iota_i,
\]
donde la primera desigualdad es lo demostrado antes y la segunda es la hipótesis de inducción.
\item Si $\lambda$ es un ordinal límite, entonces
\[
f_i^\lambda=\bigvee\{f_i^\alpha\mid \alpha<\lambda\}, \quad f^\lambda_F=\bigvee\{f^\alpha_F\mid \alpha<\lambda\}
\]
y
\[  
\iota_i\circ f_i^\lambda=\iota_i\circ \bigvee_{\alpha<\lambda}f_i^\alpha\leq\bigvee_{\alpha<\lambda}\iota_i\circ f_i^\alpha.
\]
Por la hipótesis de inducción tenemos que
\[
\iota_i\circ f_i^\alpha\leq f_F^\alpha\circ \iota_i\Rightarrow \bigvee_{\alpha<\lambda}\iota_i\circ f_i^\alpha\leq \bigvee_{\alpha<\lambda} f_F^\alpha\circ \iota_i.
\]
\end{itemize}
Así, evaluando en $0\in A_i$ se cumple que
\[
(\iota_i\circ f_i^\alpha)(0)\leq (f^\alpha_F\circ \iota_i)(0)\Rightarrow \iota_1(d_1)\leq f^\alpha_F(0\oplus 1)\;\mbox{ y }\;\iota_2(d_2)\leq f^\alpha_F(1\oplus 0).
\]
Luego 
\[
d_1\oplus d_2=d_1\oplus 1\wedge 1\oplus d_2\leq f^\alpha_F(0\oplus 1)\wedge f^\alpha_F(1\oplus 0)=f^\alpha_F(0\oplus 0). 
\]
De aquí que, para $d=d(\alpha)=f^\alpha_F(0\oplus 0)$, se cumple que
\[
(x_1\oplus x_2)\vee d_1\oplus d_2\leq (x_1\oplus x_2)\vee d,
\]
pero $x_1\oplus x_2\vee d_1\oplus d_2=(x_1\vee d_1)\oplus (x_2\vee d_2)=1\oplus 1$. Por lo tanto 
\[
x_1\oplus x_2\vee d=1\oplus 1,
\]
es decir, $A_1\oplus A_2$ es eficiente.
\end{proof}

\begin{cor}\label{Paso 3}
Si $A_1, A_2\in \Frm$ son marcos $\alpha$-eficiente y $\beta$-eficiente, respectivamente, entonces $A_1\oplus A_2$ es eficiente.
\end{cor}

\begin{proof}
Sin perdida de generalidad, consideremos $\gamma=\max\{\alpha, \beta\}$. Entonces $A_1$ y $A_2$ son $\gamma$-eficientes y por la Proposición \ref{tidyplus}, $A_1\oplus A_2$ es eficiente.
\end{proof}

Generalizamos el resultado anterior a coproductos arbitrarios.
\begin{prop}\label{coprodE}
La clase de marcos eficientes es cerrada bajo coproductos.
\end{prop}

\begin{proof}
Sean $\{A_i\}_{i\in I}$ una familia de marcos eficientes y $\{\alpha_i\}_{i\in I}\subseteq \Ord$ los grados de eficiencia de cada $A_i$. Sabemos que 
\[
\bigoplus_{i\in I} A_i\in \Frm \quad \mbox{ y }\quad (\iota_i\colon A_i\to \bigoplus_{i\in I} A_i)\in \Frm,
\]
entonces $(\iota_i)_*$ es el adjunto derecho de $\iota_i$. Además, para $F\in (\bigoplus_{i\in I} A_i)^\wedge$, por la Proposición \ref{fF}, $(\iota_i)_*F\in A_i^\wedge$ y por (\ref{Imagenfiltros}), 
\[
x_i\in (\iota_i)_*F\Leftrightarrow \iota_i(x_i)\in F.
\]
Sin perdida de generalidad, consideremos $\alpha=\sup\{\alpha_i\mid i\in I\}$, entonces cada $A_i$ es $\alpha$-eficiente. Por hipótesis, para $x_i\in (\iota_i)_*F$ se cumple que $x_i\vee d_i=1$, donde
\[
d_i=d_i(\alpha)=f_i^\alpha(0), \mbox{ con } f_i=\dot{\bigvee}\{v_y\mid y\in (\iota_i)_*F\}.
\]

\noindent
\textbf{Afirmación:} Para cada $x_i\in (\iota_i)_*F$ se cumple que $\langle x_i\rangle=\bigoplus \iota_i(x_i)\in F$.\\

Sabemos que $\bigoplus \iota_i(x_i)=\bigvee \iota_i(x_i)$. Luego, $\iota_i(x_i)\leq \bigvee \iota_j(x_j)$ para todo $i\in I$ y $j\in I$. 
Por lo tanto, al ser $F$ filtro, se cumple que $\bigvee \iota_i(x_I)\in F$.\\

Notemos que para cada $A_i$ tenemos el siguiente diagrama
\[\begin{tikzcd}
    {A_i} && {A_i} \\
    {\bigoplus_{i\in I} A_i} && {\bigoplus_{i\in I} A_i}
    \arrow["{f_i}", from=1-1, to=1-3]
    \arrow["{\iota_i}"', from=1-1, to=2-1]
    \arrow["{\iota_i}", from=1-3, to=2-3]
    \arrow["f_F"', from=2-1, to=2-3]
\end{tikzcd}\]
donde $f_F=\dot{\bigvee}\{v_{\iota_i(y)}\mid \iota_i(y)\in F\}$.\\

Si tomamos $a\in A_i$, de manera similar que en la prueba de la Proposición \ref{tidyplus}, se puede verificar que $\iota_i\circ f_i\leq f_F\circ \iota_i$ y por inducción transfinita, $\iota_i\circ f_i^\alpha\leq f_F^\alpha\circ \iota_i$. Evaluando en $0\in A_i$ se cumple que
\[
(\iota_i\circ f_i^\alpha)(0)\leq (f_F^\alpha\circ \iota_i)(0)\Leftrightarrow \iota_i(d_i)\leq f_F^\alpha(\langle 0\rangle=d(\alpha)=d.
\]
y $\langle d_i\rangle =\bigoplus \iota_i(d_i)\leq d$.\\

Luego, $\langle x_i\rangle \vee \langle d_i\rangle\leq \langle x_i\rangle \vee d$, pero 
\[
\langle x_i\rangle \vee \langle d_i\rangle=\langle x_i\vee d_i\rangle = \langle 1\rangle,
\]
es decir, $\langle 1\rangle \leq \langle x_i\rangle \vee d$ y por lo tanto $\langle x_i\rangle \vee d= \langle 1\rangle$. Así, por la afirmación anterior, $ \bigoplus_{i\in I} A_i$ es eficiente.
\end{proof} 

\section{Cuestiones en espacios topológicos}\label{Topologia}
Recordamos brevemente algunas nociones topológicas que serán utilizadas para motivar las construcciones libres de puntos posteriores.\\

Para un espacio $S$ que es $T_2$ se satisface que si $Q\in \mathcal{Q}S$, entonces $Q\in \mathcal{C}S$, donde $\mathcal{C}S$ es la colección
de subconjuntos cerrados de $S$. Podemos tener la condición anterior si tomamos una clase más general de espacios.

\begin{dfn}\label{empaquetado}
Consideremos $S\in \Top$ y $\mathcal{Q}S$ la colección de subconjuntos compactos saturados. Decimos que $S$ es \emph{empaquetado}
si para todo $Q\in \mathcal{Q}S$, entonces $Q\in \mathcal{C}S$.
\end{dfn}

De esta manera se cumple que
\[
T_2\Rightarrow \text{empaquetado}\Rightarrow T_1.
\]

Sexton relaciona los espacios empaquetados con la construcción introducida por Hochster en \cite{hochster1969prime}, conocida como el espacio de parches. Adicionalmente, 
ella define las nociones análogas a espacio de parches y espacios empaquetados, pero en marcos. Estas son el marco de parches y la trivialidad del parche mencionadas en la sección anterior.\\

Para poder tener una interacción entre los conceptos topológicos y sus análogos en marcos, necesitamos conocer las herramientas 
que nos permitan dar este salto.

%\subsection{El Teorema de Hofmann-Mislove}

\subsection{Operadores espaciales}
Debido a que $\mathcal{O}S$ es un marco, resulta natural el preguntarnos si podemos definir operadores 
$\varphi\colon \mathcal{OS}\to \mathcal{O}S$ tales que cumplan la Definición \ref{nucleo}. De manera análogo como lo hicimos en la sección 
anterior, tenemos que $N\mathcal{O}S$ es el conjunto de todos los núcleos sobre $\mathcal{O}S$. De hecho, existe una clase particular de núcleos 
de gran utilidad dentro del marco de abiertos.
\begin{dfn}\label{Nespacial}
Sean $S\in \Top$, $\mathcal{O}S\in \Frm$ y $E\subseteq S$. Para $U\in \mathcal{O}S$ definimos
\begin{equation}\label{nucleo espacial}
[E](U)=(E\cup U)^\circ.
\end{equation}
Si un operador en $\mathcal{O}S$ tiene la forma (\ref{nucleo espacial}), entonces decimos que es un \emph{núcleo espacialmente inducido}.
\end{dfn}

No es complicado verificar que (\ref{nucleo espacial}) es efectivamente un núcleo. De manera particular, 
recordado que la implicación en $\Top$ es $(W\succ U)=(W'\cup U)^\circ$ para cada $U, W\in \mathcal{O}S$ tenemos que 
\[
u_W=[W]\quad y \quad v_W=[W'].
\] 

Si a un núcleo espacialmente inducido le calculamos su complemento dual, obtenemos un operador sobre $\mathcal{C}S$. Dicho operador resulta ser una \emph{derivada idempotente}.

\begin{dfn}\label{derivadaidempotente}
Sea $S\in \Top$ y $\mathcal{C}S$ la colección de subconjuntos cerrados de $S$. Un operador $\partial\colon \mathcal{C}S\to \mathcal{C}S$ es una derivada si satisface:
\begin{enumerate}
\item Si $X\subseteq Y$ entonces $\partial(X)\subseteq \partial(Y)$,
\item $\partial(X)\subseteq X$,
\item $\partial(X\cup Y)=\partial(X)\cup \partial(Y)$
\end{enumerate}
para todo $X,Y\in \mathcal{C}S$. De manera adicional, decimos que la derivada es idempotente si $\partial^2=\partial$.
\end{dfn}

Si tomamos el núcleo $v_W=[W']$ y calculamos su complemento dual obtenemos
\begin{equation}\label{eq:derivada}
(v_{W}(U))'=((W'\cup U)^{\circ})'=\overline{(W\cap U')}=\partial_{W}(U').
\end{equation}
En este caso, la última igualdad proporciona el operador $\partial_{W}$ y es una derivada idempotente.\\

Considerando $F\in \mathcal{O}S^\wedge$, el análogo del prenúcleo (\ref{Prenucleo}) tiene la forma
\begin{equation}\label{Prenucleoespacial}
f_F=\dot{\bigvee}\{[U']\mid U\in F\}=\dot{\bigvee}\{[U']\mid Q\subseteq U\}
\end{equation}
donde la segunda igualdad se obtiene por el Teorema de Hofmann-Mislove para $Q\in \mathcal{Q}S$.\\ 

Calculando el complemento dual de (\ref{Prenucleoespacial}) obtenemos $\partial_F(X)=\bigcap\{\overline{(X\cap U)}\mid Q\subseteq U\}$
donde $X\in \mathcal{C}S$. Luego, 
\begin{equation}\label{eq:complementodual}
\partial_F(X)=f_F(X')'.
\end{equation} 
Al operador $\partial_F$ se le conoce como la $Q$-derivada (ver \cite[Lema 6.1]{simmons2004vietoris}).\\

Similarmente a como construimos la sucesión (\ref{Sucesionf}) tenemos la sucesión decreciente de $Q$-derivadas
\[
\partial_F^0=\id, \quad \partial^{\alpha+1}_F=\partial_F(\partial_F^\alpha), \quad \partial_F^\lambda=\bigcap\{\partial_F^\alpha\mid \alpha<\lambda\}.
\]

La manera de relacionar ambas sucesiones se muestra en el siguiente lema.

\begin{lem}\label{inducciontransfinita}
Para todo $\alpha\in \Ord$ y $X\in CS$, se cumple
\begin{equation}\label{eq:induccion}
\partial_F^\alpha(X)=\big(f_F^\alpha(X')\big)'.
\end{equation}
\end{lem}

\begin{proof}
Procedemos por inducción transfinita sobre $\alpha\in \Ord$.

\smallskip
\noindent\textbf{Caso base $\alpha=0$.}
Por definición, $\partial_F^0=\id$, es decir, $\partial_F^0(X)=X$ y $f_F^0=\id$, luego
\[
\partial_F^0(X)=X=(X'')=\big(f_F^0(X')\big)'.
\]

\smallskip
\noindent\textbf{Caso sucesor.}
Supongamos que la afirmación es cierta para $\alpha$ y probemos para $\alpha+1$.
Por definición de iteración,
\[
\partial_F^{\alpha+1}(X)=\partial_F\big(\partial_F^\alpha(X)\big).
\]
Además, por (\ref{eq:complementodual})
\[
\partial_F^{\alpha+1}(X)
=(f_F((\partial_F^\alpha(X))'))'.
\]
Por hipótesis inducción, $\partial_F^\alpha(X)=\big(f_F^\alpha(X')\big)'$ y sustituyendo obtenemos
\[
\partial_F^{\alpha+1}(X)=(f_F(f_F^\alpha(X')))'=\big(f_F^{\alpha+1}(X')\big)'
\]
que es lo que queríamos.

\smallskip
\noindent\textbf{Caso límite.}
Sea $\lambda$ un ordinal límite y supongamos que lo anterior se cumple para todo $\alpha<\lambda$.
Por definición de la iteración,
\[
\partial_F^\lambda(X)=\bigcap_{\alpha<\lambda}\partial_F^\alpha(X).
\]
Por De Morgan,
\[
\big(\partial_F^\lambda(X)\big)'
=\Big(\bigcap_{\alpha<\lambda}\partial_F^\alpha(X)\Big)'
=\bigcup_{\alpha<\lambda}\big(\partial_F^\alpha(X)\big)'.
\]
Usando la hipótesis de inducción, para todo $\alpha<\lambda$,
\[
\big(\partial_F^\alpha(X)\big)'=f_F^\alpha(X'),
\]
Luego,
\[
\big(\partial_F^\lambda(X)\big)'=\bigcup_{\alpha<\lambda} f_F^\alpha(X')=f_F^\lambda(X').
\]
Por lo tanto,
\[
\big(\partial_F^\lambda(X)\big)'=f_F^\lambda(X'),
\]
es decir,
\[
\partial_F^\lambda(X)=\big(f_F^\lambda(X')\big)'.
\]
y obtenemos lo que queríamos.
\end{proof}

Un caso particular del lema anterior es cuando $\alpha=\infty$.
\begin{cor}\label{cor:infinito}
Para $X\in \mathcal{O}S$ y $U\in \mathcal{O}S$ se cumple que
\begin{equation}\label{eq:infinito}
    \begin{split}
\partial_F^\infty(X)=(v_F(X'))' & \iff \partial_F^\infty(U')'=(v_F(U))\\
& \iff \partial_F^\infty(X)'=(v_F(X')).
    \end{split}
\end{equation}
\end{cor}

Agregar en esta parte el análogo a la sucesión de $d$-elementos.

\subsection{Espacios Apilados}
Si consideramos $S\in \Top$, por el Torema de Hofmann-Mislove, tenemos una correspondencia biyectiva entre 
$F\in \mathcal{O}S^\wedge$ y $Q\in \mathcal{Q}S$. Debido a que trabajaremos con filtros 
abiertos en $A$ y en $\mathcal{O}S$, necesitamos hacer diferencias entre ellos.\\ 

Denotaremos por $F$ a los filtros definidos en $A$ 
y por $\nabla$ a los filtros definidos en $\mathcal{O}S$.
Con esto en mente tenemos que los operadores espaciales asociados al filtro $\nabla\in \mathcal{O}S^\wedge$ son
\[
v_\nabla\quad\text{y}\quad \partial_\nabla^\infty. 
\]
De manera adicional, en el caso espacial, el núcleo (\ref{wF}) es $w_\nabla=[M']$. Ya que $\nabla\in \mathcal{O}S^\wedge$, se cumple que 
\begin{equation}\label{intervaloespacial}
    v_\nabla\leq [Q']\leq  w_\nabla=[M']
\end{equation}
y cuando $S$ es un espacio $T_1$ se cumple que $[Q']=[M']$.\\

Los espacios \emph{apilados} y los espacios \emph{fuertemente apilados}, además de tener relación con la eficiencia, también proporciana información
con los elementos del intervalo (\ref{intervaloespacial}).

\begin{dfn}\label{def:apilado}
Consideremos $S\in\Top$ y $\nabla\in \mathcal{O}S^\wedge$. Decimos que:
\begin{enumerate}
  \item un conjunto $Q\in \mathcal{Q}S$ es apilado si $\overline{Q}=\partial_\nabla^\infty(S)$ y es fuertemente apilado si $v_\nabla=[Q']$.
  \item el espacio $S$ es apilado o fuertemente apilado si cada $Q\in \mathcal{Q}S$ es apilado o fuertemente apilado, respectivamente.
\end{enumerate}
\end{dfn}

Notemos que si $S$ es un espacio $T_1$ y fuertemente apilado, entonces $v_\nabla=w_\nabla$, lo cual no siempre ocurre pues los espacios de Alexandroff son fuertemente apilado, pero 
no $T_1$. Algunas propiedades adicionales que satisfacen esta clase de espacios son las siguientes.

\begin{prop}\label{Propiedadesapilado}
Para $S\in \Top$ se cumple que:
\begin{enumerate}
\item Si $S$ es $T_2$, entonces $S$ es fuertemente apilado.
\item Si $S$ es fuertemente apilado, entonces $S$ es apilado.
\item Si $S$ es apilado y $T_1$, entonces $S$ es sobrio.
\end{enumerate}
\end{prop}

En la siguiente sección estudiamos la interacción entre fuertemente apilado y $T_1$. Para terminar 
esta sección, dejamos el siguiente resumen de implicaciones que tenemos en el caso espacial.
\[\begin{tikzcd}
  {\text{Eficiente}} & {\text{F. apilado}} & {\text{Apilado}} \\
  & {\text{Empaquetado}} & {T_1}
  \arrow[Rightarrow, from=1-1, to=1-2]
  \arrow["{+ T_0}"', Leftrightarrow, from=1-1, to=2-2]
  \arrow[Rightarrow, from=1-2, to=1-3]
  \arrow[Rightarrow, from=2-2, to=2-3]
\end{tikzcd}\]

Parte de las preguntas que tenemos con respecto a los espacios f. apilados y apilados son si estos tienen su análogo en marcos y, de ser así, 
¿las implicaciones anteriores se siguen cumpliendo en $\Frm$? Más adelante veremos las respuestas de estas.

\section{Interacciones entre Frm y Top}\label{Interacciones}
En esta sección introducimos versiones libres de puntos de las nociones espaciales de apilado y fuertemente apilado, y analizamos su interacción 
con la eficiencia y con la construcción de cocientes compactos. Los resultados presentados aquí forman parte del trabajo desarrollado en esta investigación.

\subsection{Empujes y levantamientos}
En esta subsección trabajamos con un caso particular de los levantamientos presentados en la Proposición \ref{levantamiento}. En este caso tenemos el diagrama
\[\begin{tikzcd}
	A & A \\
	{\mathcal{O}S} & {\mathcal{O}S}
	\arrow["j", from=1-1, to=1-2]
	\arrow["{\mathcal{U}}"', from=1-1, to=2-1]
	\arrow["k"', from=2-1, to=2-2]
	\arrow["{{\mathcal{U}}_*}"', from=2-2, to=1-2]
\end{tikzcd}\]
donde $j\in NA$, $k\in N\mathcal{O}S$ y $\mathcal{U}_*$ es el adjunto derecho de la reflexión espacial.\\

Sabemos que para $S\in \Top$, si $E\subseteq S$, entonces $[E]\in N\mathcal{O}S$. De esta manera, si $A=\mathcal{O}S$, al realizar un levantamiento a través de $\mathcal{U}$ tenemos lo siguiente.
\begin{lem}\label{lem:auxiliar}
Si $A=\mathcal{O}S$ y $E\subseteq S$ se cumple que
\[
\mathcal{U}_*[E]\mathcal{U}=[E].
\]
\end{lem}

\begin{proof}
Como el marco $A$ es espacial se cumple que $\mathcal{U}$ es un isomorfismo. De aquí que si $U\in \mathcal{O}\pt A$ existe $V\in \pt A$ tal que $U=\mathcal{U}(V)$ y como $U\in \pt A$, se cumple que $U=V$. Luego
\[
\begin{split}
(\mathcal{U}_*[E]\mathcal{U})(U)&=\mathcal{U}_*[E](\mathcal{U}(U))=\mathcal{U}_*([E](U))\\
&=\mathcal{U}_*(E\cup U)^{\circ}=\bigcup\{W\in \mathcal{O}S\mid W\subseteq (E\cup U)^{\circ}\}\\
&=(E\cup U)^{\circ}=[E](U).
\end{split}
\]
Por lo tanto $\mathcal{U}_*[E]\mathcal{U}=[E]$.
\end{proof}

En el caso general, cuando $A\in \Frm$ y $S=\pt A$, el levantamiento de un núcleo espacialmente inducido nos proporciona $j={\mathcal{U}}_*[E]\mathcal{U}\in NA$.\\

Si consideramos $F\in A^\wedge$ y $\nabla\in \mathcal{O}S^\wedge$, cada filtro abierto produce su propio intervalo de admisibilidad. Esto es
\[
[v_F, w_F]\subseteq NA\quad\text{y}\quad[v_\nabla, w_\nabla]\subseteq N\mathcal{O}S.
\]

Relacionamos ambos intervalos de admisibilidad por medio del levantamiento.

\begin{prop}\label{morfismo}
    Para $F\in A^\wedge$, $Q\in\mathcal{Q}S$ y $\nabla=\nabla(Q)$, si $k\in [v_\nabla, w_\nabla]$, entonces $\mathcal{U}_*\circ k\circ\mathcal{U}\in [v_F, w_F]$.
\end{prop}

\begin{proof}
Si $k\in [v_\nabla, w_\nabla]$, entonces $\mathcal{U}_*\circ k\circ \mathcal{U}\in NA$. Además, para $F\in A^\wedge$ si $j\in [v_F, w_F]$, entonces $F=\nabla(j)$. De esta manera debemos probar que $F=\nabla(\mathcal{U}_*\circ k\circ \mathcal{U})$.\\

Por el Teorema de Hofmann-Mislove tenemos que
\[
x\in F\iff Q\subseteq \mathcal{U}(x)\iff \mathcal{U}(x)\in \nabla=\nabla(Q).
\]

Además, para todo $x\in A$
\[
(\mathcal{U}_*\circ k\circ\mathcal{U})(x)=\mathcal{U}_*(k(\mathcal{U}(x)))=\bigwedge(S\setminus k(\mathcal{U}(x))).
\]

Como $k\in [v_\nabla, w_\nabla]$, entonces $\nabla=\nabla(k)$, es decir, para $V\in \nabla$ se cumple que $k(V)=S$. Por lo tanto,
\[
\begin{split}
x\in F\iff \mathcal{U}(x)\in \nabla &\iff S\setminus k(\mathcal{U}(x))=\emptyset\\
&\iff \bigwedge(S\setminus k(\mathcal{U}(x)))=1\\
&\iff x\in \nabla(\mathcal{U}_*\circ k\circ \mathcal{U})
\end{split}
\]
\end{proof}

Lo anterior define un morfismo $\mho\colon [v_\nabla, w_\nabla]\to [v_F, w_F]$. Podría resultar interesante analizar que clase de morfismo es $\mho$. 
Por ejemplo, notemos que ocurre con los extremos del intervalo en $NA$.\\

Consideremos $k\in [v_\nabla, w_\nabla]$. Ya que $v_F$ es un núcleo ajustado, entonces se cumple que 
\begin{equation}\label{comparacionvF}
v_F\leq \mathcal{U}_*\circ k\circ \mathcal{U}. 
\end{equation}

De manera similar, dado que $w_F$ es el mayor elemento de su bloque, entonces
\begin{equation}\label{comparacionwF}
\mathcal{U}_*\circ k\circ \mathcal{U}\leq w_F.
\end{equation}

\begin{prop}\label{igualdadwF}
Para $F\in A^\wedge$ y $\nabla\in \mathcal{O}S^\wedge$ se cumple que 
\[
w_F=\mathcal{U}_*\circ w_\nabla\circ \mathcal{U}.
\]
\end{prop}

\begin{proof}
Por (\ref{comparacionwF}), debemos verificar que $w_F\leq \mathcal{U}_*\circ w_\nabla\circ \mathcal{U}$. Por la fórmula de la adjunción,
lo anterior es equivalente a 
\[
\mathcal{U}\circ w_F\leq w_\nabla\circ \mathcal{U},
\]
es decir, para todo $x\in A$ se cumple que 
\[
\mathcal{U}(w_F(x))\subseteq [M'](\mathcal{U}(x))=(M'\cup \mathcal{U}(x))^\circ.
\]
Ya que $\mathcal{U}(w_F(x))\in \mathcal{O}S$, es suficiente demostrar que $\mathcal{U}(w_F(x))\subseteq M'\cup \mathcal{U}(x)$.
A manera de contradicción, consideremos $y\in \mathcal{U}(w_F(x))$ tal que $y\notin M'$ y $y\notin \mathcal{U}(x)$. De aqui que
\[
y\in M\quad\text{ y }\quad x\leq y.
\]

Por hipótesis, $y\in \mathcal{U}(w_F(x))$, entonces 
\[
w_F(x)\nleq y\quad \text{ y }\quad w_F(x)=\bigwedge\{m\in M\mid x\leq m\},
\]
es decir, $w_F(x)\leq y$ lo cual es una contradicción.\\ 

Por lo tanto, $\mathcal{U}(w_F(x))\subseteq M'\cup \mathcal{U}(x)$.
\end{proof}

Un caso particular del diagrama (\ref{Cuadrado}) lo podemos obtener al considerar $j^*=\mathcal{U}_A$. Así,
\[\begin{tikzcd}
	A & {A_F} \\
	{\mathcal{O}S} & {\mathcal{O}S_\nabla}
	\arrow["{v_F}", from=1-1, to=1-2]
	\arrow["{\mathcal{U}_A}"', from=1-1, to=2-1]
	\arrow["h", from=1-2, to=2-2]
	\arrow["{v_\nabla}"', from=2-1, to=2-2]
\end{tikzcd}\]

Se puede probar que $\pt A_F=Q$, es decir, tenemos la reflexión espacial
\[
\mathcal{U}_{A_F}\colon A_F\to \mathcal{O}\pt A_F=\mathcal{O}Q
\]
y este es un cociente de $A$ a través de $v_F$.\\ 

Adicionalmente, tenemos el morfismo de marcos
\[\begin{tikzcd}
	A & {\mathcal{O}S} & {\mathcal{O}S_\nabla} & {\mathcal{O}Q}
	\arrow["{\mathcal{U}_A}", from=1-1, to=1-2]
	\arrow["{v_\nabla}", from=1-2, to=1-3]
	\arrow["{\mathcal{U}_{\mathcal{O}S_\nabla}}", from=1-3, to=1-4]
\end{tikzcd}\]
y es un cociente de $A$ a través de $v_\nabla$, pues se cumple que $\pt \mathcal{O}S_\nabla=Q$. Debido a que $Q\subseteq S$, ambos cocientes están dados por la asignación
\[
x\mapsto \mathcal{U}_A(x)\cap Q
\]
para $x\in A$. 

La construcción culmina en el siguiente diagrama, que sintetiza la interacción entre cocientes compactos, separación y levantamientos de núcleos.
\begin{equation}\label{Qcuadrado}
\begin{tikzcd}
	A && {A_F} & \\
	&&& {\mathcal{O}Q} \\
	{\mathcal{O}S} && {\mathcal{O}S_\nabla}
	\arrow["{v_F}", from=1-1, to=1-3]
	\arrow["{\mathcal{U}_A}"', from=1-1, to=3-1]
	\arrow["{\mathcal{U}_{A_F}}", from=1-3, to=2-4]
	\arrow["h"', from=1-3, to=3-3]
	\arrow["{v_\nabla}"', from=3-1, to=3-3]
	\arrow["{\mathcal{U}_{\mathcal{O}S_\nabla}}"', from=3-3, to=2-4]
\end{tikzcd}\end{equation}

Dado que el diagrama anterior esta codificado a través de $F\in A^\wedge$ y el $Q\in \mathcal{Q}S$ asociado, 
denominamos a (\ref{Qcuadrado}) como el \emph{$Q$-cuadrado}.

\begin{thm}\label{haus1}
Bajo las hipótesis del $Q$-cuadrado (\ref{Qcuadrado}). Si $A\in \Frm$ satisface la propiedad $\mathbf{(H)}$, entonces
\[
\mathcal{O}Q\cong \uparrow Q'.
\]
En otras palabras, el marco de abiertos del espacio de puntos de $A_F$ es isomorfo a un cociente compacto
y cerrado de un espacio Huasdorff.
\end{thm}

\begin{proof}
Por hipótesis, $A$ satisface $\mathbf{(H)}$. Como la propiedad $\mathbf{(H)}$ es cerrada bajo cocientes, $A_F$ y $\mathcal{O}S$ satisfacen $\mathbf{(H)}$.
Además, $\mathbf{(H)}$ es conservativa, es decir, $S=\pt A$ es un espacio $T_2$. Al ser $S$ $T_2$, si $Q\in \mathcal{Q}S$, entonces $Q\in \mathcal{C}S$ y, por la Proposición \ref{Propiedadesapilado}, $S$ 
es fuertemente apilado. Por la Definición \ref{def:apilado}, $v_\nabla=[Q']$.\\

Ya que $Q'\in \mathcal{O}S$, se cumple que $[Q']=u_{Q'}$ y este es complementado en $N\mathcal{O}S$. Además,
\[
\mathcal{O}S_{u_{Q'}}=\uparrow Q'=\mathcal{O}S_{v_\nabla}.
\]

Por \cite[Propocisión VI.3.3]{PicadoPultr} se cumple que $\mathcal{O}S_{v_\nabla}$ es espacial. De aquí que $g$ es un isomorfismo.\\

Por lo tanto, $\mathcal{O}\pt A_F\cong \mathcal{O}S_\nabla$, es decir, $\mathcal{O}Q\cong \uparrow Q'$.

\end{proof}

\subsection{Nociones de apilamiento}\label{N. Apilamiento}
Proponemos a continuación las nociones libres de puntos de apilado y fuertemente apilado en la categoría $\Frm$, inspiradas en las
definiciones espaciales previas.\\

Recordemos que por la Definición \ref{def:apilado} tenemos las igualdades
\[
[Q']=v_\nabla\quad \text{y}\quad\overline{Q}=\partial_\nabla^\infty(S).
\]
que nos indican cuando un espacio $S$ es fuertemente apilado o apilado, respectivamente.\\

Si queremos tener las nociones análogas de apilado y fuertemente apilado en $\Frm$ lo ideal sería trasladar las igualdades anteriores
a través del levantamiento de $[Q']$ hacía el núcleo $v_F$. Por (\ref{comparacionvF}) se cumple que 
\[
v_F\leq \mathcal{U}_*\circ [Q']\circ \mathcal{U}.
\]

Motivados por la caracterización espacial en términos de compactos saturados y filtros en $\mathcal{O}S$, formulamos las siguientes definiciones
en la categoría $\Frm$. 
\begin{dfn}\label{def:apilado-marcos}
Sea $A\in \Frm$.
\begin{enumerate}
  \item Diremos que $A$ es fuertemente apilado si $\mathcal{U}_*[Q']\mathcal{U}=v_F$.
  \item Diremos que $A$ es apilado $(\mathcal{U}_*([Q'])\mathcal{U})(0)=v_F(0)$.
\end{enumerate}
\end{dfn}

Trivialmente tenemos que fuertemente apilado implica apilado. Nuestro objetivo es demostrar que las propiedades anteriores son conservativas, es decir, que un espacio $S$ es apilado (fuertemente apilado) si y sólo si su marco de abiertos $\mathcal{O}S$ es apilado (fuertemente apilado). 
Antes de hacer esto necesitamos el siguiente lema auxiliar.

\begin{prop}\label{apiladoconservativa}
Las nociones de fuertemente apilado y apilado son conservativas.
\end{prop}

\begin{proof}
Consideremos $S\in \Top$, $A=\mathcal{O}S$, $F\in A^\wedge$ y $Q\in \mathcal{Q}S$.\\

Primero, supongamos que $S$ es fuertemente apilado. De aquí que $v_F=[Q']$ y por el Lema \ref{lem:auxiliar} tenemos que $\mathcal{U}_*[Q']\mathcal{U}=[Q']=v_F$. Por lo tanto $\mathcal{O}S$ es fuertemente apilado.\\

Supongamos ahora que $S$ es apilado. Por lo anterior se cumple que 
\[
\overline{Q}=\partial_F^\infty(S)\iff \overline{Q}'=(\partial_F^{\infty}(S))'.
\]

Además, 
\[
\begin{split}
(\mathcal{U}_*[Q']\mathcal{U})(\emptyset)&=\mathcal{U}_*([Q'](\{V\in \pt\mathcal{O}S\mid \emptyset \not\subseteq V\}))\\
&=\mathcal{U}_*([Q'](\emptyset))=\mathcal{U}_*(Q')^{\circ}\\
&=\bigcup\{W\in \mathcal{O}S\mid W\subseteq (Q')^\circ\}\\
&=(Q')^{\circ}=\overline{Q}'.
\end{split}
\]
Por el Corolario \ref{cor:infinito} tenemos que $v_F(\emptyset)=(\partial_F^\infty(\emptyset'))'$. Por lo tanto 
\[
v_F(\emptyset)=(\mathcal{U}_*[Q']\mathcal{U})(\emptyset),
\]
es decir, $\mathcal{O}S$ es apilado.\\
\end{proof}

El análisis desarrollado en las subsecciones anteriores sugiere una posible conexión más profunda entre eficiencia y las nociones 
libres de puntos de apilado introducidas. En particular, surge la pregunta de si estas nociones permiten caracterizar la eficiencia
bajo hipótesis de separación adecuadas.\\

Las siguientes cuestiones forman parte del trabajo en curso y orientan las etapas posteriores de la investigación.

Recordemos que si un espacio es $T_1$ y fuertemente apilado, entonces se satisface que $w_\nabla=v_\nabla$. Adicionalmente, por la Proposición \ref{Propiedadesapilado}, se cumple que
si $S$ es apilado y $T_1$, entonces es sobrio.\\

Lo anterior nos hace cuestionarnos si las nociones en marcos tienen el mismo comportamiento, es decir, si $A$ es $T_1$ y fuertemente apilado, 
\[
\text{¿}v_F=w_F?
\]
Si lo anterior ocurre, ¿qué información nos proporciona con respecto al marco $A$? Además, se cumple que si $A$ es 
apilado y $T_1$, ¿$A$ es espacial?\\

Con respecto a la eficiencia, en espacios se satisface que 
\[
\text{Eficiente}\iff \text{Empaquetado y apilado}.
\]

Ya que la noción de empaquetado no es una propiedad espacial, debemos preguntarnos si existe una propiedad en marcos tal que
\[
\text{Eficiente}\iff \mathcal{P}\text{ y Apilado},
\]
donde la propiedad $\mathcal{P}$ es la noción análoga a empaquetado.

\subsection{El marco [0,1]}

Siguiendo la idea presentada en \cite{simmons2006regularity} (o, de manera alternativa, en \cite{simmons2004vietoris}), consideremos el marco linealmente ordenado 
\[
A=[0,1].
\] 
En este caso se cumplen las siguientes observaciones.
\begin{enumerate}
\item Para todo $a,b\in A$ la implicación $(a\succ b)$ está dada por
\[
(a\succ b)=\begin{cases}
1, & \text{si } a\leq b,\\
b, & \text{si } a>b.
\end{cases}
\]
\item Para $a, x\in A$, los núcleos distinguidos $u_a, v_a, w_a$ están dados por
\[
u_a(x)=\begin{cases}
  x & \text{si }a\leq x,\\
  a & \text{si }x< a,
\end{cases}, \quad v_a(x)=\begin{cases}
1 & \text{si } x\geq a,\\
a & \text{si } x<a,
\end{cases} \quad w_a(x)=\begin{cases}
1 & \text{si } x>a,\\
a & \text{si } x\leq a.
\end{cases}
\]

\item Si $F\subseteq A$ es un filtro, entonces $F$ es siempre admisible y tiene la forma
\[
F=\nabla(w_a)=(a, 1]\quad \mbox{ o }\quad F=\nabla(v_a)=[a, 1]
\]
para algún $a\in A$. En particular, $F\in A^\wedge$ si y solo si $F=(a, 1]$.

\item Si $F\in A^\wedge$, el núcleo ajustado asociado $v_F$ se calcula como
\begin{equation}\label{eq:lm}
v_F(x)=l_m(x)=\begin{cases}
1 & \text{si } x> m,\\
x & \text{si } x\leq m.
\end{cases}
\end{equation}
donde $l_m=v_m\wedge w_m$ y $m=\bigwedge F$.
\item Para $S=\pt A$ se cumple que $S=[0,1)$ y los abiertos en $\mathcal{O}S$ son de la forma
\[
\mathcal{U}_A(x)=\{p\in S\mid p<x\}=[0,x).
\]
\item El marco $A$ no es $T_1$, pues ningún $p\in S$ es máximo. En consecuencia, $A$ no satisface ninguna propiedad de separación más fuerte en $\Frm$.
\item Para $F\in A^\wedge$ el conjunto compacto saturado asociado $Q\in \mathcal{Q}S$ es de la forma $Q=[0,m]$.
\item Se cumple que $[Q']=v_F$; por lo tanto, $S$ es fuertemente apilado.
\item Ya que $[0, x)\in \mathcal{O}S$, entonces $[x, 1]\in \mathcal{C}S$ y $Q=[0, m]$. Por lo tanto, los subconjuntos compactos no son cerrados; es decir, $S$ no es empaquetado.
\item El marco $A$ no es eficiente.

\end{enumerate}

Todo lo anterior resume las principales características del marco $[0,1]$. En lo que sigue, realizaremos el análisis de las construcciones de parches tanto de 
$A$ como de su espacio de puntos $S$.\\

Recordemos que para $A\in \Frm$ y $S\in \Top$ tenemos 
\[
\text{Pbase}=\{u_a\wedge v_F\mid a\in A, F\in A^\wedge\}\quad \text{y}\quad \text{pbase}=\{U\cap Q'\mid U\in \mathcal{O}S, Q\in \mathcal{Q}S\}
\]
donde Pbase genera al marco de parches de $A$ y pbase genera la topología del espacio de parches de $S$. Por \eqref{eq:lm} tenemos que 
\[
\text{Pbase}=\{u_a\wedge l_m\mid a\in [0,1], m\in [0,1)\}.
\]

Para $x\in A$  se cumple
\[
(u_a\wedge l_m)(x)=u_a(x)\wedge l_m(x)..
\]

Analizamos por casos la relación entre $x$ y $m$.
\begin{description}
\item[Caso 1] Si $x<m$, entonces $(u_a\wedge l_m)(x)=u_a(x)\wedge x=x$.
\item[Caso 2] Si $x=m$, entonces $(u_a\wedge l_m)(x)=u_a(m)\wedge m=x$.
\item[Caso 3] Si $x>m$, entonces $(u_a\wedge l_m)(x)=u_a(x)\wedge 1=u_a(x)$.
\end{description}

En los primeros 2 casos el valor de $a$ no afecta la evaluación. La dependencia esencial aparece cuando Además $u_a\wedge l_m=\id$, es decir, los generadores básicos de la 
$m<x$, donde interviene la relación entre $a$ y $m$.
\begin{description}
\item[Caso 3.1] Si $m<x<a$, entonces $(u_a\wedge l_m)(x)=a$.
\item[Caso 3.2] Si $m<a\leq x$, entonces $(u_a\wedge l_m)(x)=x$.
\end{description}

De este análisis se deduce que, $u_a\wedge l_m=\id$ si y solo si $m<a$. Por lo tanto, lo generadores no triviales de Pbase corresponden exactamente a los pares
$(m, a)$ con $m<a$.\\

En consecuencia, el marco $PA$ queda codificado por los intervalos abiertos de la forma $(m,a)$, donde $m\in \pt A$ y $a\in A$.

Queremos identificar si en nuestro ejemplo, $PA$ satisface alguna propiedad de separación. Para hacerlo, ocuparemos recordar información respecto a los parches.\\

Si $A\in\Frm$ es un marco arbitrario y $S$ su espacio de puntos, podemos construir el siguiente diagrama conmutativo que relaciona ambas construcciones de parches.
\[\begin{tikzcd}
	A & PA & NA \\
	{\mathcal{O}S} & {P\mathcal{O}S} & {N\mathcal{O}S} \\
	& {\mathcal{O}^pS} & {\mathcal{O}^fS}
	\arrow[from=1-1, to=1-2]
	\arrow["{\mathcal{U}_A}"', from=1-1, to=2-1]
	\arrow[hook, from=1-2, to=1-3]
	\arrow["{P(\mathcal{U}_A)}", from=1-2, to=2-2]
	\arrow["{N(\mathcal{U}_A)}", from=1-3, to=2-3]
	\arrow[from=2-1, to=2-2]
	\arrow[hook, from=2-1, to=3-2]
	\arrow[hook, from=2-2, to=2-3]
	\arrow["\pi", from=2-2, to=3-2]
	\arrow["\sigma", from=2-3, to=3-3]
	\arrow[hook, from=3-2, to=3-3]
\end{tikzcd}\]
donde $\mathcal{O}^fS$ es la topología de Skula (o topología frontal) del espacio $S$. En \cite{sexton2006point} lo llaman el 
\emph{diagrama del ensamble de parches} y más detalles sobre este también pueden ser encontrados
en \cite{sexton2003point}.\\

En nuestro caso, $\mathcal{U}_A$ es un isomorfismo. Además, bajo $\mathcal{U}_A$, la construcción de parches resulta ser funtorial y, por lo tanto,
$PA\cong P\mathcal{O}S$. De manera similar, $NA\cong N\mathcal{O}S$.
Para describir explícitamente $P\mathcal{O}S$ y $N\mathcal{O}S$ utilizamos el siguiente resultado, que puede consultarse en \cite{AvilaBezhaniMorandiZaldivar2020}.

\begin{prop}\label{Ensamble espacial}
Sea $S$ un conjunto parcialmente ordenado.
\begin{enumerate}
  \item Si $S$ no tiene anticadenas infinitas, entonces $N\mathcal{O}S$ es espacial.
  \item Si $S$ es totalmente ordenado, entonces $N\mathcal{O}S$ es espacial.
\end{enumerate}
\end{prop}

Así, por la Proposición \ref{Ensamble espacial} tenemos que $N\mathcal{O}S$ es espacial y por lo tanto $N\mathcal{O}S\cong \mathcal{O}^fS$. Juntando toda esta información 
el diagrama del ensamble de parches en nuestro caso particular queda como sigue.

\[\begin{tikzcd}
	A & PA & NA \\
	{\mathcal{O}S} & {\mathcal{O}^pS} & {\mathcal{O}^fS}
	\arrow[from=1-1, to=1-2]
	\arrow["\cong"', from=1-1, to=2-1]
	\arrow[from=1-2, to=1-3]
	\arrow["\cong", from=1-2, to=2-2]
	\arrow["\cong", from=1-3, to=2-3]
	\arrow[from=2-1, to=2-2]
	\arrow[from=2-2, to=2-3]
\end{tikzcd}\]
Por lo tanto, si queremos estudiar al marco $PA$ podemos estudiar a la topología de parches $\mathcal{O}^pS$. Recordemos que 
\[
\text{pbase}=\{U\cap Q'\mid U\in \mathcal{O}S, Q\in \mathcal{Q}S\}.
\]
En nuestro caso,
\[
\text{pbase}=\{[0, a)\cap [m, 1)\mid a\in [0,1], m\in [0,1)\}=\{(m,a)\}.
\]

Como la topología generada por los intervalos $(m,a)$ coincide con la topología usual en $[0,1)$, el espacio $^pS$ es Hausdorff. En consecuencia, $\mathcal{O}^pS\cong PA$ satisface $\mathbf{(H)}$, aun cuando $A$ no es $T_1$ ni eficiente.

\bibliographystyle{amsalpha}
\bibliography{research2}

\end{document}